\chapter{Marco teórico}\label{cap:marco}

\section{Lavado de dinero en criptomonedas}

El lavado de dinero es el proceso por el cual se oculta el origen ilícito de unos fondos para
integrarlos en la economía formal. La literatura lo describe en tres etapas: colocación, en la
que los fondos entran al sistema; estratificación, en la que se mueven a través de muchas
operaciones para borrar su rastro; e integración, en la que vuelven a la economía como fondos
aparentemente legítimos \cite{Chen2018AMLReview}. Las criptomonedas ofrecen a esas etapas un
medio rápido, transfronterizo y seudónimo, y han dado lugar a técnicas propias como la mezcla
de fondos, el encadenamiento de transacciones y el uso de plataformas de intercambio con
controles débiles \cite{Almeida2025Techniques, Subashi2024Crypto}.

Los sistemas tradicionales de monitoreo se basan en reglas fijas, como umbrales de monto o
listas de contrapartes, y producen una proporción muy alta de falsas alarmas
\cite{Kute2021AMLReview}. El aprendizaje automático permite aprender patrones a partir de casos
etiquetados \cite{Chen2018AMLReview, Vassallo2021Boosting, Pocher2023Anomalous}. En Bitcoin, la
cadena de bloques es pública y las transacciones están enlazadas por los fondos que se
transfieren, de modo que el problema se puede plantear como la clasificación de los nodos de
un grafo \cite{Weber2019Elliptic}.

\section{Redes neuronales de grafos}\label{sec:gnn}

Sea un grafo con $N$ nodos, matriz de atributos $\mathbf{X} \in \mathbb{R}^{N \times d}$ y
conjunto de vecinos $\mathcal{N}(v)$ de cada nodo $v$. Una red neuronal de grafos calcula en
cada capa $\ell$ una representación $\mathbf{h}_v^{(\ell)}$ de cada nodo combinando su propia
representación con las de sus vecinos \cite{Scarselli2009GNN, Wu2021Survey, Zhou2020Review}:
\begin{equation}
  \mathbf{h}_v^{(\ell+1)} = \phi\!\left(\mathbf{h}_v^{(\ell)},\;
  \bigoplus_{u \in \mathcal{N}(v)} \psi\!\left(\mathbf{h}_v^{(\ell)}, \mathbf{h}_u^{(\ell)}\right)\right),
  \qquad \mathbf{h}_v^{(0)} = \mathbf{x}_v,
\end{equation}
donde $\bigoplus$ es una agregación invariante al orden (suma, media o máximo) y $\phi$ y
$\psi$ son funciones con parámetros aprendidos. Tras $L$ capas, la representación de un nodo
depende solo de los nodos a $L$ saltos o menos, su subgrafo receptivo. La última capa produce
la probabilidad de cada clase. Las cuatro arquitecturas del estudio difieren en cómo combinan
el nodo con sus vecinos.

\textbf{GCN} \cite{Kipf2017GCN}. Promedia al nodo con sus vecinos usando la matriz de adyacencia
normalizada por el grado, con un único conjunto de pesos:
\begin{equation}
  \mathbf{h}_v^{(\ell+1)} = \sigma\!\left(\mathbf{W}^{(\ell)} \sum_{u \in \mathcal{N}(v) \cup \{v\}}
  \frac{\mathbf{h}_u^{(\ell)}}{\sqrt{(d_u + 1)(d_v + 1)}}\right).
\end{equation}
El nodo entra en la suma como un vecino más y no tiene un peso propio.

\textbf{GraphSAGE} \cite{Hamilton2017GraphSAGE}. Separa la representación del nodo de la de sus
vecinos con dos matrices de pesos:
\begin{equation}
  \mathbf{h}_v^{(\ell+1)} = \sigma\!\left(\mathbf{W}_1^{(\ell)} \mathbf{h}_v^{(\ell)} +
  \mathbf{W}_2^{(\ell)} \frac{1}{|\mathcal{N}(v)|} \sum_{u \in \mathcal{N}(v)} \mathbf{h}_u^{(\ell)}\right).
\end{equation}

\textbf{GAT} \cite{Velickovic2018GAT}. Pondera a cada vecino, incluido el propio nodo, con un
coeficiente de atención aprendido:
\begin{equation}
  \mathbf{h}_v^{(\ell+1)} = \sigma\!\left(\sum_{u \in \mathcal{N}(v) \cup \{v\}} \alpha_{vu}
  \mathbf{W}^{(\ell)} \mathbf{h}_u^{(\ell)}\right), \qquad
  \alpha_{vu} = \operatorname{softmax}_u\!\left(\operatorname{LeakyReLU}\!\left(\mathbf{a}^\top
  [\mathbf{W}\mathbf{h}_v \,\|\, \mathbf{W}\mathbf{h}_u]\right)\right),
\end{equation}
y repite el cálculo con varias cabezas de atención independientes cuyos resultados se
concatenan.

\textbf{TAGCN} \cite{Du2017TAGCN}. Aplica un filtro polinomial de orden $K$ sobre la matriz de
adyacencia normalizada $\tilde{\mathbf{A}}$, con una matriz de pesos por cada potencia:
\begin{equation}
  \mathbf{H}^{(\ell+1)} = \sigma\!\left(\sum_{k=0}^{K} \tilde{\mathbf{A}}^{k} \mathbf{H}^{(\ell)}
  \mathbf{W}_k^{(\ell)}\right).
\end{equation}
El término $k = 0$ da al nodo un peso propio, y cada capa alcanza vecinos a $K$ saltos.

\section{Desbalance de clases}\label{sec:desbalance}

Cuando una clase es mucho más frecuente que otra, la pérdida promedio está dominada por la
clase mayoritaria y el modelo puede minimizarla prediciendo casi siempre esa clase
\cite{He2009Imbalanced, Johnson2019Survey}. Las estrategias para compensarlo actúan sobre los
datos o sobre la pérdida.

Sobre los datos, se puede submuestrear la clase mayoritaria o sobremuestrear la minoritaria.
SMOTE \cite{Chawla2002SMOTE} genera ejemplos sintéticos de la clase minoritaria interpolando
entre un ejemplo y uno de sus vecinos más cercanos de la misma clase:
$\tilde{\mathbf{x}} = \mathbf{x}_i + \lambda(\mathbf{x}_j - \mathbf{x}_i)$, con
$\lambda \sim \mathcal{U}(0, 1)$.

Sobre la pérdida, sea $p_i$ la probabilidad que el modelo asigna a la clase verdadera del
ejemplo $i$. La entropía cruzada sin ajuste es $\mathcal{L} = -\frac{1}{n}\sum_i \log p_i$.
Con pesos por clase, cada ejemplo se pondera por el inverso de la frecuencia de su clase,
$w_c = n / (2 n_c)$, de modo que
\begin{equation}
  \mathcal{L}_{\text{PC}} = \frac{1}{2}\,\frac{1}{n_{+}}\sum_{i:\,y_i = 1} (-\log p_i) +
  \frac{1}{2}\,\frac{1}{n_{-}}\sum_{i:\,y_i = 0} (-\log p_i),
\end{equation}
y cada clase aporta la mitad de la pérdida sin importar cuántos ejemplos tenga. La
\textit{focal loss} \cite{Lin2020Focal} reduce el peso de los ejemplos que el modelo ya
clasifica bien:
\begin{equation}
  \mathcal{L}_{\text{FL}} = -\frac{1}{n}\sum_i \alpha_{y_i} (1 - p_i)^{\gamma} \log p_i,
\end{equation}
donde $\gamma$ controla cuánto se atenúan los ejemplos fáciles y $\alpha_{y_i}$ pondera la
clase. Con $\alpha = 0{,}75$ para la clase rara y $0{,}25$ para la frecuente, un ejemplo
ilícito pesa tres veces lo que uno negativo, mientras que con pesos por clase y una razón de
41 a 1 pesa 41 veces más.

\textbf{Métricas bajo desbalance.} La exactitud no sirve: con un 2~\% de ilícitas, predecir
siempre la clase negativa da un 98~\% de aciertos. La curva ROC relaciona la tasa de
verdaderos positivos con la de falsos positivos, y como esta última se divide por el número
enorme de negativos, puede verse alta aunque las falsas alarmas superen con creces a los
aciertos. La curva de precisión y exhaustividad sí refleja ese costo, y su área (PR-AUC) tiene
como referencia del azar la prevalencia de la clase positiva \cite{Saito2015PR}. El
coeficiente de correlación de Matthews resume la matriz de confusión completa y es más
informativo que el F1 bajo desbalance \cite{Chicco2020MCC}.

\section{Explicabilidad de modelos}\label{sec:xai}

La explicabilidad busca que las decisiones de un modelo puedan ser comprendidas por una persona
\cite{Adadi2018Peeking, Guidotti2018Survey, Arrieta2020XAI}. Los métodos post hoc explican un
modelo ya entrenado sin modificarlo, y los locales explican una predicción individual
asignando una importancia a cada elemento de la entrada. En grafos, esos elementos pueden ser
las variables de los nodos, las aristas o subgrafos completos \cite{Yuan2023Taxonomic,
Dai2024Trustworthy}. Esta tesis estudia explicaciones de variables: para una transacción, cuál
de sus 165 variables pesó más en la predicción.

\textbf{GNNExplainer} \cite{Ying2019GNNExplainer}. Busca una máscara $\mathbf{M}$ sobre las
variables (y opcionalmente sobre las aristas) que preserve la predicción del modelo y sea a la
vez pequeña:
\begin{equation}
  \min_{\mathbf{M}} \; -\log P\!\left(\hat{y} \mid \mathbf{X} \odot \sigma(\mathbf{M})\right)
  + \lambda_1 \lVert \sigma(\mathbf{M}) \rVert_1 + \lambda_2 \, H\!\left(\sigma(\mathbf{M})\right),
\end{equation}
donde $H$ es la entropía de la máscara, que la empuja hacia valores cercanos a 0 o 1. La
máscara se inicializa al azar y se optimiza por descenso de gradiente, por lo que dos
ejecuciones pueden terminar en máscaras distintas.

\textbf{PGExplainer} \cite{Luo2020PGExplainer}. Entrena una red que, a partir de las
representaciones de los dos extremos de una arista, predice su importancia. Una vez entrenada,
explica cualquier nodo sin optimizar de nuevo. Atribuye importancia a aristas, no a variables.

\textbf{Valores de Shapley} \cite{Lundberg2017SHAP, Strumbelj2014Explaining}. Tratan a las
variables como jugadores de un juego cooperativo cuyo pago es la predicción. La importancia de
la variable $j$ es su contribución marginal promedio sobre todos los órdenes posibles en que
las variables se incorporan:
\begin{equation}
  \varphi_j = \frac{1}{|\Pi|}\sum_{\pi \in \Pi} \left[f\!\left(S_{\pi}^{j} \cup \{j\}\right)
  - f\!\left(S_{\pi}^{j}\right)\right],
\end{equation}
donde $S_{\pi}^{j}$ es el conjunto de variables que preceden a $j$ en el orden $\pi$ y las
variables ausentes toman un valor de referencia. Como hay $d!$ órdenes, se estima con una
muestra aleatoria de permutaciones \cite{Castro2009Polynomial}, y por ello es estocástico.

\textbf{Integrated Gradients} \cite{Sundararajan2017IG}. Integra el gradiente de la salida a lo
largo del camino recto entre una referencia $\mathbf{x}'$ y la entrada $\mathbf{x}$:
\begin{equation}
  \mathrm{IG}_j(\mathbf{x}) = (x_j - x'_j) \int_0^1 \frac{\partial f\!\left(\mathbf{x}' +
  \alpha(\mathbf{x} - \mathbf{x}')\right)}{\partial x_j}\, d\alpha.
\end{equation}
La integral se aproxima con una suma de Riemann. Con una referencia fija, el método es
determinista. Si la referencia es el vector cero y una variable vale cero, su atribución es
exactamente cero.

\textbf{Expected Gradients} \cite{Erion2021Expected}. Sustituye la referencia fija por una
distribución de referencias $D$ tomadas de los datos y promedia:
$\mathrm{EG}_j(\mathbf{x}) = \mathbb{E}_{\mathbf{x}' \sim D,\, \alpha \sim \mathcal{U}(0,1)}
\left[(x_j - x'_j)\, \partial_j f(\mathbf{x}' + \alpha(\mathbf{x} - \mathbf{x}'))\right]$.
La esperanza se estima por muestreo, por lo que es estocástico.

\section{Evaluación de las explicaciones}

No existe una única medida de la calidad de una explicación. Las revisiones sistemáticas
distinguen varias propiedades \cite{Nauta2023Anecdotal, Bodria2023Benchmarking}: la
\textbf{fidelidad}, si la explicación refleja lo que de verdad usa el modelo; la
\textbf{plausibilidad}, si coincide con lo que un experto considera relevante; y la
\textbf{estabilidad} o consistencia, si explicaciones de la misma predicción en condiciones
equivalentes se parecen entre sí. En grafos, las evaluaciones sistemáticas han mostrado que
los explicadores difieren mucho entre sí y que su desempeño depende del conjunto de datos y de
la arquitectura \cite{Agarwal2023Evaluating, Longa2024Explaining}.

La estabilidad es una condición necesaria para usar una explicación como evidencia: si dos
ejecuciones del mismo explicador sobre el mismo modelo y la misma transacción dan órdenes de
importancia distintos, el analista no sabe a cuál creer. Pero no es suficiente. Un explicador
que ignorara el modelo y devolviera siempre el mismo orden sería perfectamente estable e
inútil. Adebayo et al. \cite{Adebayo2018Sanity} propusieron por eso comparar la explicación de
un modelo entrenado con la de la misma arquitectura con pesos aleatorios: si se parecen, la
explicación no depende de lo aprendido.

Para medir el acuerdo entre dos explicaciones se usan comparaciones de rangos o de conjuntos.
La correlación de Spearman \cite{Spearman1904} compara los órdenes completos de importancia, y
el índice de Jaccard \cite{Jaccard1912} compara los conjuntos de las $k$ variables más
importantes. También se han propuesto métricas diseñadas específicamente para la estabilidad
de las explicaciones \cite{Gawantka2024Explanations}.
